\SourceLeaf{22}
\[
[1]\qquad \xi=\int\!\int\frac{\mathrm{d}p.\mathrm{d}q.\cos p.\sin p}{Ar}
\]
\par
Ebenso erhalten wir nach dem sechsten Lehrsatz
\[
\begin{aligned}
[2]\qquad \xi&=-\int\!\int\frac{\mathrm{d}p.\mathrm{d}q.\sin p}{r^3}.(a-x)\bigl(\frac{(a-x)x}{AA}+\frac{(b-y)y}{BB}+\frac{(c-z)z}{CC}\bigr)
\end{aligned}
\]
\par
Schließlich liefert uns der vierte Lehrsatz
\[
\begin{aligned}
[3]\qquad &\int\!\int\frac{\mathrm{d}p.\mathrm{d}q.\sin p}{r^3}\bigl(\frac{(a-x)x}{AA}+\frac{(b-y)y}{BB}+\frac{(c-z)z}{CC}\bigr)=0\\
&\hspace{6em}\text{oder }=-\frac{4\pi}{ABC}
\end{aligned}
\]
je nachdem, ob der Punkt $M$ außerhalb oder innerhalb des Körpers liegt.

Wir betrachten nun die Größen $A,B,C$ als besondere Werte dreier Variablen $\alpha,\beta,\gamma$, die folgendermaßen gewählt sind. Die Differenzen $\alpha\alpha-\beta\beta$ und $\alpha\alpha-\gamma\gamma$ sollen konstant sein. So lässt sich $\xi$ als Funktion der Variablen $\alpha,\beta,\gamma$ oder vielmehr einer von ihnen ansehen. Die gleichzeitigen Änderungen der Größen $\xi,\alpha,\beta,\gamma$ bezeichnen wir durch das Zeichen $\delta$. Aus Gleichung $[1]$ folgt leicht, dass bei unbegrenztem Wachstum von $\alpha,\beta,\gamma$ die Größe $\xi$ über alle Grenzen hinaus abnimmt, da offenbar selbst der kleinste Wert von $r$ über alle Grenzen hinaus wächst. Wir müssen also $\xi=0$ für $\alpha=\infty$ setzen. Differenziert man Gleichung $[1]$ in der Form
\[
\alpha\xi=\int\!\int\frac{\mathrm{d}p.\mathrm{d}q.\cos p.\sin p}{r}
\]
mit dem Zeichen $\delta$ für die Variation, so ergibt sich
\[
\alpha\delta\xi+\xi\delta\alpha=-\int\!\int\frac{\mathrm{d}p.\mathrm{d}q.\cos p.\sin p.\delta r}{rr}
\]
Nun gilt
\[
\begin{aligned}
r\delta r&=-(a-x)\delta x-(b-y)\delta y-(c-z)\delta z\\
&=-(a-x)\cos p.\delta\alpha-(b-y)\sin p.\cos q.\delta\beta-(c-z)\sin p\sin q.\delta\gamma\\
&=-(a-x)x.\frac{\delta\alpha}{\alpha}-(b-y)y.\frac{\delta\beta}{\beta}-(c-z)z.\frac{\delta\gamma}{\gamma}\\
&=-\alpha\delta\alpha.\bigl(\frac{(a-x)x}{\alpha\alpha}+\frac{(b-y)y}{\beta\beta}+\frac{(c-z)z}{\gamma\gamma}\bigr)
\end{aligned}
\]
wegen $\alpha\delta\alpha-\beta\delta\beta=0$, $\alpha\delta\alpha-\gamma\delta\gamma=0$. Daraus ergibt sich
\[
\alpha\delta\xi+\xi\delta\alpha=\delta\alpha.\int\!\int\frac{\mathrm{d}p.\mathrm{d}q.x\sin p}{r^3}\bigl(\frac{(a-x)x}{\alpha\alpha}+\frac{(b-y)y}{\beta\beta}+\frac{(c-z)z}{\gamma\gamma}\bigr)
\]
\par
Zieht man hiervon Gleichung $[2]$, mit $\delta\alpha$ multipliziert, ab, nachdem man $A,B,C$ durch $\alpha,\beta,\gamma$ ersetzt hat, so ergibt sich
